% Nouveau EGA V (cote 25): archive cote 25, site page 58 (PDF page 59).
% typescript. Draft compared: none.
% AI-transcribed by Claude Opus 5.5 (claude-opus-5-5, adaptive thinking, effort high), Message Batches API, batch ega25b (msgbatch_01A7n5YuXubUPNK1zsepi8RG), 2026-11-13; checked against both supplied scan parts; see CHECK_DECISIONS.tsv.
\scanpage{25/58}
\mspage{53}

10. \uline{Pinceaux de sections hyperplanes}, \uline{et fibrations de variétés éclatées}

10. 1. Soit Z l'ensemble P-exceptionnel dans $\overset{\hand{\vee}}{P}$ relativement à une
certaine propriété constructible\inserted{P}, de sorte que Z est une partie construc%
tible de $\overset{\hand{\vee}}{P}$ . \inserted[typed]{Supposons S=Spec(k),} On sait (cf nº\hand{12}\deleted{suivant}, où nous \unsure{faisans|faisons} un ratrappage
% NOTE: a long diagonal pencil stroke in the left margin beside this and the next two lines
\margin{$/$}\inserted{de choses} qui aurai\over{t}{\unsure{ent}} sans doute dû venir dans les nºs antérieurs) que pour que
l'on ait codim\deleted{\unsure{et}}(Z,$\overset{\hand{\vee}}{P}$ ) $\hand{\geqslant}$ 2, il f et s que \hand{``}\xout{t}toute droite assez
générale" L dans $\overset{\hand{\vee}}{P}$ ne rencontre pas Z, \xout{et il} ou encore $\overset{\hand{\scriptscriptstyle -}}{Z}$ , - et il suffit
qu'il existe \uline{une} droite L dans $\overset{\hand{\vee}}{P}$ ne rencontrant pas k. Lorsque \xout{\unsure{e|c}tte con}
k est infini, il f donc et s qu'il existe une doite L \margin{\unsure{$-$|$\wedge$}}dans $\overset{\hand{\vee}}{P}$ ne rencon-%
trant pas $\overset{\hand{\scriptscriptstyle -}}{Z}$ . On appelle "pinceau linéaire\deleted{"}\inserted[typed]{de sections hyperplanes de X\hand{''}} défini par la\xout{d}droite
L de $\overset{\hand{\vee}}{P}$\hand{,} le L-prés\unsure{c|f}héma $\underline{Y}_L$ (définition valable pour un S quel\xout{f}conque).
Ainsi, les réflexions précédentes, jointes aux résultats des nºs 8 et 9,
nous donnent des critères d'existence de tels pinceaux, ayant des fibres
$\underline{Y}_{\hand{\unsure{\xi|\zeta}}}$ ( \hand{$\unsure{\xi|\zeta} \in$}L) satisfaisant toutes la propriété P, tout d'abord dans le
cas \xout{\unsure{d'un|deux}} \inserted[typed]{où S est un} corps de base infini. \xout{Mais précédant comme} Compte tenu de
8.2.\hand{,} lorsque \xout{X} pour tout cycle premier associé à X, on a dim f(\unsure{F|T})$\hand{>}$0,
on peut (en prenant la propriété P' = P + condition de régularité
pour $\hand{\unsure{\phi|\emptyset}}_{\hand{\unsure{\xi|\zeta}}}$ ) exiger que le pinceau $\underline{Y}_L$ soit plat sur L. Dans le cas
où S est quelconque, on peut encore\xout{x} souvent, par le procédé de 7.1. ,
construire un tel pinceau au dessus d'un voisinage ouvert d'un point
donné s de S, pourvu que k(s) soit infini et qu'on sache que Z est fermé
(ce qui sera assuré dans divers cas par les résultats du par.5 et
l'hypothèse de propreté pour X $\hand{\rightarrow}$ S).

\margin{|} Pour bien faire, il conviendrait, après des explications géné-%
rales de cette eau, de donner de\unsure{\xout{x}s|s\xout{x}} énoncés récapitulatifs où on
applique effectivement les résultats précédents pour un certain nombre
de propriétés de cette nature (y compris bien sur des propriétés modulées\hand{)}

\mspage{58}
